\section{Methodology}
\label{sec:methodology}

\subsection{Background interpolation around a blinded mass window}

At each mass hypothesis, the search must predict a smooth trident background without
allowing a possible narrow resonance to influence that prediction. We therefore fit the
GP on a fixed, dataset-specific mass range and exclude a small region around the test
mass from training. The GP interpolates across the excluded region, and those same bins
are then used in the signal-plus-background likelihood. Earlier HPS studies redefined
the sidebands in separate piecewise windows. The fixed-support construction gives
smoother behavior across the scan and allows the three datasets to enter a joint
likelihood with a shared coupling.

\subsection{Search interval and GP fit support}
\label{sec:fit-support-policy}

The search interval specifies where a resonance is tested; the wider GP fit
support specifies which bins can constrain its background. At each mass, bins
inside the moving $\pm2.25\sigma_m$ exclusion are omitted from GP training.
Keeping training data beyond the search endpoints avoids unnecessarily
one-sided interpolation. The current settings are given in
Table~\ref{tab:v4p9-card-lineage} and apply to both standalone and combined fits.
The 2015 extension holds the 90~MeV kernel coordinates fixed above 90~MeV;
Appendix~\ref{app:v504-extension} gives its numerical and conditional checks.

\begin{table}[htbp]
\centering
\begin{tabular}{lrrr}
\toprule
Data & Search [MeV] & GP support [MeV] & Upper length factor \\
\midrule
2015 full & 19--100 & 14--135 & 8 \\
2016 full & 39--180 & 30--210 & 12 \\
2021 10\% & 50--250 & 36--300 & 15 \\
\bottomrule
\end{tabular}
\caption{Search ranges and GP settings used for the individual and combined fits.
The 2021 lower edge follows the support study;
its earlier 40~MeV value is retained only in historical comparisons.}
\label{tab:v4p9-card-lineage}
\end{table}

The 2016 lower edge at 30~MeV also reflects the actual bin geometry. At the
39~MeV search endpoint, the blind interval begins at 35.063~MeV, whereas a
35~MeV support starts at a bin center of 35.125~MeV and supplies no lower-side
training center. Training masks are defined using bin centers, not only
continuous interval endpoints. The 2021 support study is described in
Section~\ref{sec:v4p9p5-support}; its post-selection practical rule and remaining
55~MeV under-recovery remain part of the qualification of these settings.

\subsection{Observable, scan grid, and blinding}

The observable is the reconstructed $e^+e^-$ invariant-mass histogram for each
dataset. Throughout this section, $d$ labels the dataset, $m$ denotes the reconstructed
$e^+e^-$ invariant mass, $m_0$ is the test-mass hypothesis, and $\sigma_d(m)$ is the
dataset-specific mass resolution introduced in
Section~\ref{sec:event-selection}. The mass-hypothesis spacing is
\begin{equation}
\Delta m = \SI{1}{MeV},
\end{equation}
and the blind window for a mass hypothesis $m_0$ is centered on the detector
resolution,
\begin{equation}
\mathcal{B}_d(m_0) =
\left[m_0 - n_{\mathrm{blind}}\sigma_d(m_0),\,
      m_0 + n_{\mathrm{blind}}\sigma_d(m_0)\right].
\end{equation}
The extraction half-width and GP training exclusion are independently configurable.
Early injection studies used
$n_{\mathrm{blind}}=1.64$ with the same $1.64\,\sigma_d$ training exclusion. The nominal
analysis uses a
wider $2.25\,\sigma_d$ blind/extraction window with the GP training exclusion also set
to $2.25\,\sigma_d$. For an ideal Gaussian line shape, the $1.64\,\sigma_d$ interval
contains
\begin{equation}
f_{\mathrm{blind}}^{\mathrm{ideal}}(1.64)
= \mathrm{erf}\!\left(\frac{1.64}{\sqrt{2}}\right)
\simeq 0.899,
\end{equation}
while the $2.25\,\sigma_d$ interval contains
\begin{equation}
f_{\mathrm{blind}}^{\mathrm{ideal}}(2.25)
= \mathrm{erf}\!\left(\frac{2.25}{\sqrt{2}}\right)
\simeq 0.976.
\end{equation}
We normalize the signal over the full Gaussian line shape and use only its blinded-bin
slice in the local likelihood; the slice is not renormalized inside the window.

The notation $\mathcal{B}_d(m_0)$ denotes the set of histogram bins that lie inside the
blinded interval for dataset $d$ at test mass $m_0$. In the nominal fit there is no
additional gap between the extraction window and the GP
training sample: when the extraction and training-exclusion half-widths are equal, the
training exclusion starts immediately outside $\mathcal{B}_d(m_0)$. A wider exclusion,
for example $1.96\,\sigma_d$ instead of $1.64\,\sigma_d$, reduces signal leakage into
the training region but also increases the interpolation distance and therefore the GP
predictive uncertainty. In the simple background-dominated proxy
$S/\sqrt{B}\propto F(k)/\sqrt{k}$ with $F(k)=\mathrm{erf}(k/\sqrt{2})$, the optimum lies
near $k\simeq1.4$, so the earlier $1.64\,\sigma_d$ extraction was close to the
pure-sensitivity optimum. Section~\ref{sec:toys} then motivates the current
$2.25\,\sigma_d$ blind/extraction choice as a signal-recovery improvement after
full-refit pseudoexperiments showed appreciable training-region absorption at narrower
exclusions.

We assign bins by their fixed centers rather than splitting bins that overlap a
boundary. Let $c_i=(e_i+e_{i+1})/2$ be the center of bin $i$. The full
bin is included in the blinded likelihood when
\begin{equation}
m_0-n_{\mathrm{blind}}\sigma_d(m_0)
\le c_i \le
m_0+n_{\mathrm{blind}}\sigma_d(m_0),
\end{equation}
and it enters the sideband sample only when its center lies strictly outside the
training-exclusion interval. The same center mask is applied to the counts, GP
prediction, covariance, and signal template, so no observed count is divided between
the two regions.

The Gaussian signal template is integrated over the full bin edges in
Eq.~\eqref{eq:window-template}, but the selected set of bins changes discretely when a
boundary passes a bin center. Consequently the quoted local significance and
$\eps^2$ limit are not meaningful to a precision finer than the mass binning and scan
grid used to define these masks. The same discrete procedure is used for the data and
pseudoexperiments.

\subsection{Log-space preprocessing and the choice \texorpdfstring{$\alpha=1/y$}{alpha=1/y}}

Following recent HEP GPR background-modeling studies
\cite{BarrLiu2025,Gandrakota2023}, we fit the GP in transformed coordinates:
\begin{equation}
x = \log m,
\qquad
z =
\begin{cases}
\log y, & y>0, \\
0, & y=0,
\end{cases}
\end{equation}
where $y$ is the observed bin count. The log transformation is used for two related
reasons. First, the prompt trident spectrum spans a wide dynamic range, and modeling
$z$ rather than $y$ turns multiplicative shape variations into approximately additive
ones. Second, $x=\log m$ makes equal distances in the GP coordinate correspond to equal
fractional changes in mass, which is the relevant scale once the detector resolution is
written in the transformed coordinate as $\sigma_x(m)\simeq \sigma_d(m)/m$. Barr and
Liu emphasize that applying logarithms to both axes improves convergence and
performance for wide, steeply falling spectra \cite{BarrLiu2025}.
The GP regression step uses scikit-learn's GaussianProcessRegressor estimator
\cite{ScikitLearn}; the surrounding HPS analysis code supplies the histogram
handling, mass-dependent sideband masks, signal templates, profile likelihood,
and \CLs{} limit construction.

The GP is therefore used to model an unobserved smooth mean function
$z_{\mathrm{true},d}(x)$ for the log-count spectrum. In that language, ``latent'' means
only that the smooth log-rate is not observed directly; it is inferred from the
histogram counts through the transformed coordinates and the heteroscedastic variance
term. The GP mean vector is the vector of latent log-count means evaluated at the
selected training-bin coordinate. For Poisson counts $Y_i \sim
\mathrm{Pois}(\lambda_i)$, the delta method gives
$\mathrm{Var}[\log Y_i]\approx \mathrm{Var}[Y_i]/E[Y_i]^2 \approx 1/\lambda_i$, which
motivates the plug-in choice $\alpha_i\approx 1/y_i$ for occupied bins. A
heteroscedastic diagonal term $\alpha_i$ is then added to the training covariance,
\begin{equation}
\alpha_i =
\begin{cases}
1/y_i, & y_i > 0, \\
1, & y_i = 0,
\end{cases}
\label{eq:alpha-oney}
\end{equation}
\begin{equation}
K_{\mathrm{obs}} = K(X,X) + \mathrm{diag}(\bm{\alpha}),
\end{equation}
so that each $\alpha_i$ plays the role of an observation-variance term for bin $i$.
This is the quantity that GaussianProcessRegressor adds to the diagonal during fitting,
so it is treated here as a variance model rather than an arbitrary tuning parameter.

Here $\alpha_i$ is the latent-space variance associated with the transformed counts;
we do not add an extra $\log y$ weight. For empty bins the floor
$\alpha_i=1$ is a unit variance in log-count space, not a pseudo-count of one event. It
simply prevents zero-count bins from receiving ill-defined or disproportionate weight
after the log transformation. In this high-statistics search such bins are rare, so the
floor acts mainly as a conservative numerical safeguard rather than as a dominant
source of uncertainty.

The distinction from Ref.~\cite{BarrLiu2025} is more specific than a difference in
kernel choice. Barr and Liu also use a Constant$\times$RBF kernel and optimize its
amplitude and length scale. Their independent methodological study uses a spectrum
derived from a published CMS search; it is not a CMS Collaboration GPR result. In
addition to optimizing the kernel, they treat the diagonal array as tunable
regularization. For occupied bins their nominal prescription is
\begin{equation}
\alpha_i^{\mathrm{BL}}
= \frac{\sqrt{y_i}}{y_i\log y_i}
= \frac{1}{\sqrt{y_i}\log y_i},
\end{equation}
and fixed regional multipliers $f_i$ are then applied to
$\alpha_i^{\mathrm{BL}}$. Reducing the first eleven multipliers to 0.1 raised the
fraction of 100 background-only pseudoexperiments passing their BumpHunter criterion
from 51\% to 87\%, but localized residual structures remained; their high-luminosity
benchmark subsequently extended the reduced-multiplier region to the first twenty
bins \cite{BarrLiu2025}. The study therefore demonstrates that diagonal tuning can
improve a particular benchmark, while also showing that the tuned region is
sample- and regime-dependent.

The active HPS configuration applies no regional multiplier to the early bins and
fixes Eq.~\eqref{eq:alpha-oney} as the delta-method variance model. This choice matters
because scikit-learn's alpha parameter has both an observation-noise interpretation
and a numerical
Tikhonov-regularization interpretation \cite{ScikitLearnGPRDocs16}. HPS adopts the
former and places its separately validated smoothness control in the allowed
length-scale domain. These are complementary controls: changing $\alpha_i$ changes
the leverage of bin $i$ on the covariance diagonal, whereas changing $\ell$ changes
the off-diagonal correlation range of the latent function. No matched study presented
here establishes that one strategy is intrinsically superior to the other.

\subsection{Kernel choice and distinct regularization roles}
\label{sec:kernel-choice}

We model the background with a single ConstantKernel$\times$RBF kernel,
\begin{equation}
k(x,x') = C \exp\!\left[-\frac{(x-x')^2}{2\ell^2}\right],
\end{equation}
with broad amplitude bounds and a resolution-scaled RBF length-scale domain. This is
the same simple kernel class used in the methodological studies of
Refs.~\cite{Gandrakota2023,BarrLiu2025}, but the surrounding noise and validation
choices need not be the same. The two quantities optimized in every HPS sideband fit
are $C$ and $\ell$. The ConstantKernel value $C=k(x,x)$ is the pointwise prior
variance of the latent log-rate, so $\sqrt{C}$ is its prior standard deviation; $C$
is not a fitted event-count normalization. The length scale $\ell$ controls how
quickly covariance falls with separation in $x=\log m$, and therefore how rapidly
the background function may bend. The fixed $\alpha_i$ values are not a third kernel
parameter and there is no optimized WhiteKernel in the baseline. They enter only
through the observation-covariance diagonal.

Figure~\ref{fig:gpr-hyperparameter-roles} separates these roles graphically. This
separation is useful when discussing regularization: changing $C$ rescales the
covariance amplitude, changing $\ell$ changes the correlation geometry and smoothness
of the latent function, and changing $\alpha_i$ changes how strongly one observation
constrains that function.

\begin{figure}[H]
\centering
\graphicorplaceholder{0.99\linewidth}{methodology_figs/gpr_hyperparameter_roles.pdf}
\caption{Roles of the three quantities in the GPR covariance. Panel (a) shows the
normalized RBF correlation for three illustrative values of
$\ell/\sigma_x(m_0)$; at a separation equal to $\ell$, every curve has correlation
$e^{-1/2}\simeq0.61$. The length scale sets the correlation range. Panel (b) varies
$C$ at fixed $\ell$. It changes the
covariance amplitude, with pointwise prior standard deviation $\sqrt{C}$, without
changing the normalized correlation range. Panel (c) holds the kernel fixed and
changes one diagonal variance $\alpha_j$, thereby changing that bin's statistical
weight. We optimize $C$ and $\ell$ on the sidebands and fix $\alpha_i$ from the
count-based variance rule. The curves are schematic.}
\label{fig:gpr-hyperparameter-roles}
\end{figure}

We use the Constant$\times$RBF model. A Constant$\times$Mat\'ern kernel with
$\nu=3/2$ and piecewise kernels remain exploratory robustness checks; they have not
been tested for coverage in the three-campaign configuration.

\subsection{Resolution-scaled length-scale parameterization and bounds}
\label{sec:length-scale}

Because the GP is trained in $x=\log m$, detector resolution must be translated into
log-mass units. We define
\begin{equation}
\sigma_x(m) = \log\!\left(m+\sigma_d(m)\right) - \log m
\simeq \frac{\sigma_d(m)}{m}
\qquad (\sigma_d \ll m),
\end{equation}
and use this local physical scale to set the RBF bounds. Within the search intervals used here, the local term sets both bounds,
\begin{equation}
\ell_{\min,d}(m) = k_{\min,d}\sigma_x(m),
\qquad
\ell_{\max,d}(m) = k_{\max,d}\sigma_x(m).
\label{eq:dataset-lengthscale-bounds}
\end{equation}
For the RBF covariance used here,
\begin{equation}
k_{\mathrm{RBF}}(x,x') =
\sigma_f^2\exp\!\left[-\frac{(x-x')^2}{2\ell^2}\right],
\qquad
\frac{k_{\mathrm{RBF}}(x,x+\ell)}{k_{\mathrm{RBF}}(x,x)}
=e^{-1/2}\simeq0.61 .
\label{eq:rbf-lengthscale-interpretation}
\end{equation}
Thus $\ell$ is a covariance-decay scale in invariant mass: larger values
favor smoother background functions, while smaller values admit more rapid
variation \cite{RasmussenWilliams2006,FrateGP2017}. The ratio
$\ell/\sigma_x(m)$ expresses that nuisance-model smoothness relative to the width of
the narrow signal hypothesis. It is not, without an independently validated mapping,
a detector resolution measurement or a physical background correlation length.
The fitted $\ell_{\mathrm{opt}}$ is the log-marginal-likelihood optimum conditional
on the kernel, data support, sideband mask, noise model, and configured optimizer
domain. Consequently an optimum at $\ell_{\max}$ establishes only that the numerical
constraint is active; it neither measures a physical scale at the boundary nor
locates the unconstrained optimum. The implementation optimizes log-transformed
kernel hyperparameters within these configured bounds
\cite{ScikitLearnGPRDocs16}.

For a fixed dataset $d$ and mass hypothesis $m_0$, let
$\bm{z}_{\mathrm{sb}}$ contain only the transformed sideband observations and define
\begin{equation}
K_y(C,\ell)
= C\,R_{\ell}(X_{\mathrm{sb}},X_{\mathrm{sb}})
+ \mathrm{diag}(\bm{\alpha}_{\mathrm{sb}}).
\end{equation}
The optimized pair maximizes
\begin{equation}
\log p(\bm{z}_{\mathrm{sb}}\mid X_{\mathrm{sb}},C,\ell)
= -\frac{1}{2}\bm{z}_{\mathrm{sb}}^{T}K_y^{-1}\bm{z}_{\mathrm{sb}}
  -\frac{1}{2}\log|K_y|
  -\frac{n_{\mathrm{sb}}}{2}\log(2\pi).
\label{eq:gpr-log-marginal-likelihood}
\end{equation}
The quadratic term rewards agreement with the sideband observations, while the
determinant term accounts for the covariance volume used to obtain that agreement.
The implementation searches in the logarithms of $C$ and $\ell$ with multiple
optimizer restarts. This is an empirical-Bayes, or type-II maximum-likelihood,
selection of the covariance parameters; it is not a profile of $\ell$ as a
frequentist nuisance parameter in the subsequent local \CLs{} likelihood.

Figure~\ref{fig:gpr-mass-hypothesis-optimization} shows why the optimum is specific
to a dataset and mass hypothesis. Moving $m_0$ changes which bins are excluded from
training and changes the resolution-scaled domain through $\sigma_x(m_0)$. The same
dataset can therefore return a different $(C_{\mathrm{opt}},\ell_{\mathrm{opt}})$ at
neighboring hypotheses even though the kernel family is unchanged. An interior
maximum means only that the configured range is nonbinding for that fit. A maximum on
a boundary means that the range is active, while a long shallow likelihood ridge means
that the parameters are weakly identified. None of these cases turns
$\ell_{\mathrm{opt}}$ into a direct measurement of detector resolution or a physical
background correlation length.

\begin{figure}[H]
\centering
\graphicorplaceholder{0.99\linewidth}{../figures/v501_kernel_optimization_clean.pdf}
\caption{Meaning of per-mass-hypothesis hyperparameter optimization. Panel (a)
shows one schematic sideband problem in
$u=(x-x_0)/\sigma_x(m_0)$. The shaded $2.25\,\sigma_x$ interval is excluded from
training; the three curves illustrate the different interpolation behavior allowed
by short, interior, and long candidate length scales. Panel (b) shows the joint
log-marginal-likelihood surface for the same synthetic sidebands. The star is the
bounded maximum in $(C,\ell)$, and the elongated contours expose their possible
tradeoff. Panel (c) shows that moving $m_0$ changes both the sideband mask and the
local resolution scale that defines the allowed $\ell$ interval. Panel (d) illustrates
how a scan records an interior optimum differently from repeated upper-bound contact:
the former is a conditional best fit within a nonbinding range, whereas the latter
diagnoses an active numerical domain and motivates a predeclared range and closure
study. The figure is explanatory only; it contains no HPS observations or fit
results.}
\label{fig:gpr-mass-hypothesis-optimization}
\end{figure}

This distinction determines how the hyperparameter study enters a physics result.
If the allowed range forces $\ell$ too low, the GP retains unnecessary
signal-scale flexibility and can absorb part of a localized excess; if the fitted
background is forced too smooth, broad mismodeling can instead leave a biased local
residual. Raising only the upper bound permits, but does not require, a smoother
solution. A larger dataset also need not prefer a larger $\ell$: additional
statistics can resolve background structure and move the optimum downward. The
upper range is therefore selected from boundary occupancy, likelihood ordering,
optimizer stability, and fixed-amplitude signal-recovery diagnostics, never from a
favorable observed limit or local $p_0$. These checks can justify a bounded candidate
range, but they do not by themselves establish frequentist coverage.

The three-campaign analysis uses the dataset-specific factors in
Table~\ref{tab:production-lengthscale-bounds}. The lower factors come from full-refit
closure studies; changing the 2016 upper factor does not requalify those studies at
full statistics.

\begin{table}[H]
\centering
\begin{tabular}{lccc}
\toprule
Dataset & $k_{\min,d}$ & $k_{\max,d}$ & Basis \\
\midrule
2015 full & 1.0 & 8 & full-refit closure study \\
2016 full & 0.9 & 12 & closure plus post-selection range study \\
2021 10\% & 1.1 & 15 & matched-refit closure study \\
\bottomrule
\end{tabular}
\caption{Resolution-scaled RBF length-scale factors used in the three-campaign
analysis. The factors multiply $\sigma_x(m)$ in
Eq.~\eqref{eq:dataset-lengthscale-bounds}. The 2016 upper factor was chosen after a
range limitation appeared in the observed scan; it is therefore a post-selection
choice, not a pre-unblinding or coverage qualification.}
\label{tab:production-lengthscale-bounds}
\end{table}

Figure~\ref{fig:v501-lengthscale-regions} displays the allowed GP length-scale
regions in log-mass coordinates. The shaded region is a configured domain, not
an error band on the fitted smoothness. The implementation retains the upper-bound
floor $0.8\,k_{\max,d}\operatorname{median}_{m\in\mathcal M_d}\sigma_x(m)$,
with the median evaluated on its configured 300-point search grid; that floor is
nonbinding over the three plotted search ranges.
\fig{0.99}{../figures/v501_lengthscale_allowed_regions.pdf}{Allowed RBF
length scales in $x=\log m$ across the three search ranges. Shading lies between
the lower and upper bounds evaluated from the frozen analysis card; the dotted
curves show the local detector-resolution scale $\sigma_x(m)$. The lower/upper
factors are $1.0/8$, $0.9/12$, and $1.1/15$ for 2015, 2016, and 2021,
respectively. These are numerical domains for the sideband fit, not fitted
length scales or detector-resolution uncertainties.}
{fig:v501-lengthscale-regions}

The lower factors are based on full-refit functional-form closure scans, with the
observed-spectrum result kept out of the selection criterion. For 2015, the 100-toy
scan slightly prefers $k_{\min}=1.1$ in its aggregate score, but $k_{\min}=1.0$ has
similar pull-width calibration and a weaker expected-reach proxy. We retain 1.0 as the
more conservative choice. For 2016, the 25-toy continuation supports the
change to $k_{\min}=0.9$: its nonzero-injection RMS pull mean is 0.321 and its median
pull width is 0.948, compared with 0.403 and 0.889 for $k_{\min}=1.1$. It gives the
smallest aggregate calibration score among the scanned 0.9, 1.0, and 1.1 rows and
reduces the coherent apparent signal-recovery bias that motivated the follow-up. For
2021, the corrected matched-reference scan finds small rank separations among the
same three settings. The $k_{\min}=1.1$ row has the median pull width closest to
unity and the best extraction-level $\eps^2$ proxy, while the 0.9 and 1.0 rows have
slightly smaller aggregate scores. We choose 1.1 as a calibration--reach tradeoff, not
because it wins every diagnostic. These studies used the data fractions stated in
Section~\ref{sec:toys} and do not establish full-statistics closure.

The initial upper factors were 8, 8, and 15 for 2015, 2016, and 2021. In the reviewed
observed scan, the 2016 optimum reached its upper bound at all 142 hypotheses. We
therefore varied only $k_{\max,2016}$, evaluating factors
8, 10, 12, 15, and 20 on the same 39--180~MeV mass grid. A fit is counted as
upper-bound occupied when $\ell_{\mathrm{opt}}/\ell_{\max}\geq0.999$. The respective
occupancies are 142, 56, 0, 0, and 0 of 142. Factor 12 is the first scanned setting
with no occupied hypotheses. It is also followed by a numerical plateau: from factor
12 to 15, the largest pointwise change in the 2016 observed yield upper limit is
0.083\%, the largest $|\Delta Z|$ is 0.00255, and the log-marginal-likelihood
differences lie between $-4.5\times10^{-6}$ and $8.0\times10^{-6}$. The factor-15 to
factor-20 comparison is similarly stable. The selection is based on removal of the
artificial boundary and the next-setting plateau, not on which observed limit is
tighter or which local $p_0$ is smaller.

Because the observed spectrum revealed the range limitation, factor 12 is a
post-selection choice for the three-campaign analysis. It is not a pre-unblinding
choice or a coverage-calibrated setting; full-statistics functional-form closure and
direct coverage remain separate questions.
The configurations continue to retain a dataset-wide median-resolution floor on
$\ell_{\max,d}$ so that the ceiling does not become artificially small where the
local resolution is exceptionally tight. This is a guard based on a representative
$\sigma_x$ value, not a floor derived from luminosity or event statistics.
Alternative upper bounds,
fixed-background fits, and the common-lower-bound overlays remain robustness studies
collected in the appendices.

An initial 2021 source/exposure study used a nested-Poisson construction. It motivated
the later matched-refit studies but is not part of the procedure used here.

Figure~\ref{fig:gpr-window-construction} summarizes the same construction in the
language used by the implemented likelihood. The sideband bins set the log-count
training sample and its heteroscedastic bin-noise assignment, the kernel hyperparameters
are optimized only on those sidebands, and the trained GP is then conditioned to predict
the mean and covariance in the blinded mass window. That predicted count-space mean and
covariance are the background model used in the local profiled signal extraction.

\begin{figure}[H]
\centering
\begin{subfigure}[t]{0.52\linewidth}
  \centering
  \graphicorplaceholder{\linewidth}{methodology_figs/gpr_training_three_step_rail.png}
  \caption{Log-count transform, noise-aware covariance training, and window prediction.
  The filled covariance boxes indicate the sideband-sideband, sideband-blind, and
  blind-blind kernel blocks used when conditioning the GP.}
\end{subfigure}
\hfill
\begin{subfigure}[t]{0.45\linewidth}
  \centering
  \graphicorplaceholder{\linewidth}{methodology_figs/blinded_mass_window_reworked.png}
  \caption{Sideband-only training and blinded-window prediction. Filled markers show
  binned count data, vertical error bars show the fixed per-bin statistical noise, and
  the horizontal reference lines are count-axis guides used to compare sideband and
  blinded-window scales.}
\end{subfigure}
\caption{Kernel and blinded-window construction. The GP is
trained on sideband bins with fixed bin-noise variances, not on the blinded bins. The
posterior prediction in the blind window supplies the background mean and covariance
that enter the local likelihood; the shaded band in the schematic denotes the GP
posterior interval and the dashed vertical lines mark the blinded-window boundaries.}
\label{fig:gpr-window-construction}
\end{figure}

\subsection{Blind-window prediction in count space}
\label{sec:blind-window-prediction}

For each mass point, we train the GP only on the bins outside the training
exclusion window and predict the posterior mean and covariance of the latent smooth
log-rate in the blinded bins. Since the fit is performed in log-count space, that
posterior is mapped back to count space through the standard lognormal moment
relations:
\begin{align}
\mu_i &= \exp\!\left(\mu^{(\log)}_i + \frac{1}{2}\Sigma^{(\log)}_{ii}\right), \\
\mathrm{Cov}(\lambda_i,\lambda_j)
&= \mu_i\mu_j\left[\exp\!\left(\Sigma^{(\log)}_{ij}\right)-1\right].
\end{align}
The resulting blind-window mean vector $\bm b_d(m)$ and latent-background
covariance $\bm C_d^{\rm GP}(m)$ supply the local extraction. Poisson variation
of the observed counts is handled separately in the likelihood. Thus this
covariance describes uncertainty in the predicted smooth mean, not the
counting variance of the histogram.

Before factorization, the numerical implementation symmetrizes the covariance
and applies the small diagonal conditioning recorded by the solver checks.
Below, $\bm C_d$ denotes that conditioned covariance, with
\begin{equation}
 \bm C_d\equiv\bm C_{d,\rm eff}=L_dL_d^\top.
 \label{eq:Cd-cholesky}
\end{equation}
The lower-triangular factor $L_d$ maps independent unit Gaussian coordinates
into correlated count-space deformations. The regularization permits stable
factorization; it is not an added physics uncertainty.

\subsection{Profiled background implementation}
\label{sec:v501-profile-implementation}

At one tested mass, write the local Poisson mean as
\begin{equation}
 \bm\lambda_d(A_d,\bm\theta_d)
 =\bm b_d+L_d\bm\theta_d+A_d\bm w_d.
 \label{eq:v501-profile-mean}
\end{equation}
Here $A_d$ is the total Gaussian signal yield and $\bm w_d$ is the
bin-integrated template restricted to the likelihood window, as defined in
Section~\ref{sec:v501-signal-conversion}. The nuisance vector $\bm\theta_d$
has a unit Gaussian constraint, so a displacement $L_d\bm\theta_d$ preserves
the correlations predicted by the GP. At fixed $A_d$, the fit minimizes
\begin{equation}
 \ell_d(A_d,\bm\theta_d)=
 \sum_i\left[\lambda_{d,i}-n_{d,i}
 +n_{d,i}\log\!\left(\frac{n_{d,i}}{\lambda_{d,i}}\right)\right]
 +\frac12\bm\theta_d^\top\bm\theta_d,
 \qquad \lambda_{d,i}>0.
 \label{v5-eq:profile-objective}
\end{equation}
The $n\log n$ contribution is zero at $n=0$. Apart from constants independent
of the fit, $\ell_d$ is the negative log likelihood. Its first term measures
the Poisson mismatch; the second penalizes background deformations in units
of the correlated GP uncertainty. The profiled objective is
\begin{align}
 \widehat{\widehat{\bm\theta}}_d(A_d)
 &=\operatorname*{arg\,min}_{\bm\theta_d}\ell_d(A_d,\bm\theta_d),\\
 \ell_{p,d}(A_d)
 &=\ell_d\!\left(A_d,\widehat{\widehat{\bm\theta}}_d(A_d)\right).
 \label{eq:v501-profile-definition}
\end{align}
For an interior solution, the nuisance gradient and Hessian are
\begin{align}
 \nabla_{\theta}\ell_d
 &=L_d^\top\!\left(\bm1-\frac{\bm n_d}{\bm\lambda_d}\right)+\bm\theta_d,\\
 H_{\theta\theta}
 &=L_d^\top\operatorname{diag}\!\left(\frac{n_{d,i}}{\lambda_{d,i}^{2}}\right)L_d+I,
 \label{eq:v501-profile-curvature}
\end{align}
where the vector division is elementwise. These expressions make the
correlated adjustment explicit: the fit changes a set of nuisance modes,
rather than choosing unrelated background shifts in each bin.

\methodfig{0.99}{../figures/v501_profile_background_flow.pdf}{At each tested signal
yield, the background deforms within the fixed GP constraint and the
minimized objective is recorded. Repeating this optimization constructs the
profile curve used for extraction and limits. The diagram describes the
calculation; it contains no fitted data.}{fig:v501-profile-flow}

During this local profile, $\bm b_d$, $L_d$, the bin selection and the signal
template remain fixed. Profiling $\bm\theta_d$ does not retrain the GP kernel.
A complete pseudoexperiment may first retrain the sideband prediction under
its declared kernel policy, then apply this same local fit. Positivity
applies to the total Poisson mean $\lambda_{d,i}$; a separate positivity
constraint is not imposed on $\bm b_d+L_d\bm\theta_d$. The production solver
uses numerical safeguards near zero means; accepted interior solutions are
checked against the objective above. Signed auxiliary signal fits describe
both excesses and deficits, while physical discovery and upper-limit tests
use a nonnegative signal.

This Gaussian-constrained construction follows the profile-likelihood
framework of Refs.~\cite{HistFactory,Cowan2011,FrateGP2017}. The constraint
comes from sideband interpolation and should not be described as an
independent auxiliary measurement. Its qualification therefore includes the
background and signal-recovery tests presented later.

\subsection{Signal model and conversion to \texorpdfstring{$\eps^2$}{epsilon2}}
\label{sec:v501-signal-conversion}

\methodfig{0.99}{../figures/v501_yield_to_coupling_flow.pdf}{From the selected mass
spectrum to a fitted signal yield or upper bound, then to the corresponding
coupling. The campaign factor contains the declared radiative-fraction and
prompt-density normalization. The applicable visible branching factor is
included when interpreting the result above the dimuon threshold.}
{fig:v501-yield-conversion}


For each dataset $d$ and mass hypothesis $m_0$, we first construct a
full-range Gaussian line-shape template over the histogram binning,
\begin{equation}
g_{d,i}(m_0) =
\int_{e_i}^{e_{i+1}}
\frac{1}{\sqrt{2\pi}\sigma_d(m_0)}
\exp\!\left[-\frac{(m-m_0)^2}{2\sigma_d(m_0)^2}\right]\,dm,
\end{equation}
where $e_i$ and $e_{i+1}$ are the lower and upper edges of histogram bin $i$.
The template is normalized on the stored histogram support:
\begin{equation}
w^{\mathrm{full}}_{d,i}(m_0)=\frac{g_{d,i}(m_0)}{\sum_{j\in\mathrm{full\ range}}g_{d,j}(m_0)},\qquad
\sum_{i \in \mathrm{full\ range}} w^{\mathrm{full}}_{d,i}(m_0) = 1.
\end{equation}
The template entering the local likelihood is then the blinded-bin slice of this
full-range object,
\begin{equation}
w_{d,i}(m_0) =
\begin{cases}
w^{\mathrm{full}}_{d,i}(m_0), & i \in \mathcal{B}_d(m_0), \\
\text{not used in the local fit}, & i \notin \mathcal{B}_d(m_0).
\end{cases}
\label{eq:window-template}
\end{equation}
By the blinded-bin slice we mean exactly this restriction of the full normalized
Gaussian template to the subset of bins in $\mathcal{B}_d(m_0)$; the template is not
renormalized after the restriction.
Its blinded-bin sum is therefore
\begin{equation}
f_{\mathrm{blind},d}(m_0) = \sum_{i \in \mathcal{B}_d(m_0)} w_{d,i}(m_0) < 1,
\end{equation}
with a value near 0.90 for a $1.64\,\sigma_d$ half-width and near 0.98 for the current
$2.25\,\sigma_d$ observed-limit half-width. The signal
amplitude $A_d$ is defined as the \emph{total} local Gaussian signal yield associated
with the mass hypothesis, not merely the yield inside the blind window. The expected
signal contribution entering bin $i$ of the local likelihood is therefore
\begin{equation}
s_{d,i}(A_d;m_0) = A_d\,w_{d,i}(m_0),
\label{eq:signal-bin-contribution}
\end{equation}
so the total expected signal actually visible inside the blind window is
$f_{\mathrm{blind},d}(m_0)\,A_d$.

This convention preserves the fact that a finite blind region never contains the
entire Gaussian line shape. It
also keeps the signal parameter used in the fit aligned with the full-histogram
injection model used in the procedural toy studies.

The conversion between signal yield and dark-photon coupling follows the standard
fixed-target relation between $A^\prime$ production and radiative trident production
\cite{EssigEtAl2011,HPS2015DarkPhoton,HPS2016PromptLong}. In the present notation,
\begin{equation}
A_d\equiv N_{A^\prime,d}
= \frac{3\pi}{2N_{\mathrm{eff}}^{\mathrm{BR}}\alpha_{\mathrm{EM}}}
\; m_{A^\prime}\,\eps_{\rm phys}^2\,
\frac{dN_{\gamma^\ast}}{dm}\bigg|_{m=m_{A^\prime}},
\end{equation}
where $\alpha_{\mathrm{EM}}$ is the fine-structure constant and $3\pi/2$ is the
narrow-width coefficient relating on-shell production to radiative tridents.
Here $\eps_{\rm phys}^2$ is the physical kinetic-mixing parameter and
$N_{\mathrm{eff}}^{\mathrm{BR}}=\Gamma_{\mathrm{tot}}/\Gamma_{ee}
=1/\mathcal B_{ee}$ includes the phase-space dependence of all open decay
channels. In the minimal model it equals one below the dimuon threshold.
For the electron-pair yield, define
\begin{equation}
 \eps_{ee}^2\equiv\frac{\eps_{\rm phys}^2}{N_{\mathrm{eff}}^{\mathrm{BR}}}.
\end{equation}
The following fit and combination equations abbreviate this electron-channel
coordinate as $\eps^2$. The displayed minimal-visible interpretation multiplies
it by $N_{\mathrm{eff}}^{\mathrm{BR}}$ above $2m_\mu$, as specified in the
dimuon-threshold subsection. Writing the
radiative differential rate as
\begin{equation}
\frac{dN_{\gamma^\ast}}{dm}
= f_{\mathrm{rad},d}(m)\,\rho_d(m),
\end{equation}
where $\rho_d(m)$ is the differential prompt background rate in the selected
invariant-mass spectrum, we convert between $A_d$ and $\eps^2$ through
\begin{equation}
\eps^2 =
\frac{2\alpha_{\mathrm{EM}}\,A_d}
{3\pi\,m\,f_{\mathrm{rad},d}(m)\,\rho_d(m)},
\label{eq:eps2-conversion}
\end{equation}
where $\rho_d(m)$ is estimated from the observed histogram as a local counts-per-GeV
density in the separately declared prompt-density interval,
\begin{equation}
\rho_d(m)\equiv
\frac{N_d\!\left[m-n_{\mathrm{blind}}\sigma_d(m),\,m+n_{\mathrm{blind}}\sigma_d(m)\right]}
{2\,n_{\mathrm{blind}}\sigma_d(m)},
\end{equation}
with $n_{\mathrm{blind}}=1.64$ retained as the current prompt-density normalization
window even when the extraction blind window is widened to $2.25\,\sigma_d$. This keeps
the yield-to-$\eps^2$ conversion tied to the established local density scale while
the likelihood itself can recover the wider signal tails. The density-window width
remains configurable so that alternative conventions can be reproduced when needed.
It is convenient to define the
dataset-specific conversion factor
\begin{equation}
K_d(m) \equiv A_d(\eps^2 = 1)
= \frac{3\pi\,m\,f_{\mathrm{rad},d}(m)\,\rho_d(m)}{2\alpha_{\mathrm{EM}}},
\label{eq:Kd-conversion}
\end{equation}
so that a shared coupling $\eps^2$ predicts a dataset-specific total signal yield
$A_d = K_d(m)\eps^2$.

Unless stated otherwise, the methodology section uses the nominal radiative fraction
$f_{\mathrm{rad},d}(m)$. The optional penalty that propagates normalization uncertainty
into $\eps^2$ is treated as a systematic variation rather than as part of the core GP
fit and is discussed separately in Section~\ref{sec:radfrac-penalty}. The numerical
penalty values are therefore result-configuration choices, not ingredients of the
background interpolation or local extraction fit. When that systematic option is
enabled for dataset $d$, the effective radiative fraction is
\begin{equation}
f_{\mathrm{rad},d}^{\mathrm{eff}}(m)
= \left(1-\delta_{f,d}\right) f_{\mathrm{rad},d}(m),
\end{equation}
so the same signal-yield limit maps to a weaker $\eps^2$ limit.

\subsection{Statistical inference in a single dataset}

For a fixed dataset $d$ and mass hypothesis $m$, let $\bm{n}_d$ denote the vector of
observed blinded-bin counts, $\bm{b}_d$ the GP-predicted blinded-bin mean,
$\bm{C}_d=L_dL_d^\top$ the corresponding covariance of
Section~\ref{sec:blind-window-prediction} and Eq.~\eqref{eq:Cd-cholesky}, and
$\bm{w}_d(m)$ the
non-renormalized blinded-bin slice of the full Gaussian template defined in
Eq.~\eqref{eq:window-template}. The profiled local likelihood is
\begin{equation}
\mathcal{L}_d(A_d,\bm{\theta}_d)
=
\left[
\prod_{i=1}^{N_d(m)}
\mathrm{Pois}\!\left(
n_{d,i}\,\middle|\,
b_{d,i} + (L_d\bm{\theta}_d)_i + A_d w_{d,i}(m)
\right)
\right]
\exp\!\left(-\frac{1}{2}\bm{\theta}_d^\top\bm{\theta}_d\right),
\label{eq:singlelike}
\end{equation}
where the Gaussian nuisance vector $\bm{\theta}_d\sim\mathcal{N}(0,I)$ induces the
correlated background-deformation field $L_d\bm{\theta}_d$ defined below
Eq.~\eqref{eq:Cd-cholesky}, and $N_d(m)$ is the number of bins in the local blinded
set $\mathcal{B}_d(m)$. The signal contribution entering the Poisson mean is exactly
the blinded-bin restriction of Eq.~\eqref{eq:signal-bin-contribution}.

At fixed $m$, profiling Eq.~\eqref{eq:singlelike} with no physical sign constraint on
$A_d$ yields the unconstrained extraction estimator $\hat A_d^{\,\mathrm{unc}}$ and
its local uncertainty $\sigma_{A,d}$. This unconstrained fit is used for closure,
linearity, and pull studies because it allows downward background fluctuations to be
captured by the estimator itself rather than being clipped at zero. For discovery and
upper limits, however, the parameter of interest is restricted to the
physical region $A_d\ge0$, and the corresponding physical best fit is
\begin{equation}
\hat A_d = \max\!\left(0,\hat A_d^{\,\mathrm{unc}}\right).
\end{equation}
We denote by $\hat{\bm{\theta}}_d$ the nuisance-parameter maximizer at the
unconditional best fit and by $\hat{\hat{\bm{\theta}}}_d(A_d)$ the conditional maximum
likelihood estimator of $\bm{\theta}_d$ for fixed $A_d$.

\subsubsection{Local significance}

The single-dataset discovery statistic is the bounded profile-likelihood ratio
\begin{equation}
q_{0,d}(m) = -2\ln
\frac{\mathcal{L}_d\!\left(A_d=0,\hat{\hat{\bm{\theta}}}_d(0)\right)}
{\mathcal{L}_d\!\left(\hat A_d,\hat{\bm{\theta}}_d\right)},
\label{eq:q0-single}
\end{equation}
with $\hat A_d\ge0$. The local one-sided discovery $p$-value and the corresponding
Gaussian significance are then obtained with the asymptotic Cowan-et-al.\ mapping
\cite{Cowan2011}. The resulting local-significance definition is explicitly
one-sided and discovery facing: downward fluctuations are allowed in the unconstrained
fit used for closure studies, but they do not generate a positive local discovery
significance. We use this profiled multi-bin test throughout.

\subsubsection{\texorpdfstring{\CLs}{CLs} upper limits}

To test a fixed non-negative signal amplitude $A_d$, the analysis uses
the bounded one-sided profile-likelihood statistic
\begin{equation}
\tilde q_{A_d}(m)=
\begin{cases}
0, &
\hat A_d^{\,\mathrm{unc}} > A_d, \\
-2\ln
\dfrac{\mathcal{L}_d\!\left(A_d,\hat{\hat{\bm{\theta}}}_d(A_d)\right)}
{\mathcal{L}_d\!\left(\hat A_d^{\,\mathrm{unc}},\hat{\bm{\theta}}_d\right)}, &
0\le \hat A_d^{\,\mathrm{unc}}\le A_d, \\
-2\ln
\dfrac{\mathcal{L}_d\!\left(A_d,\hat{\hat{\bm{\theta}}}_d(A_d)\right)}
{\mathcal{L}_d\!\left(0,\hat{\hat{\bm{\theta}}}_d(0)\right)}, &
\hat A_d^{\,\mathrm{unc}} < 0,
\end{cases}
\label{eq:qmu-single}
\end{equation}
which is the standard bounded Cowan construction for upper limits \cite{Cowan2011}.
For a positive observed statistic, the corresponding modified-frequentist tail areas are
\begin{equation}
\mathrm{CL}_{s+b}(A_d;m)
=
P\!\left(
\tilde q_{A_d}\ge \tilde q_{A_d}^{\mathrm{obs}}
\middle|\,
A_d+\mathrm{bkg}
\right),
\qquad
\mathrm{CL}_{b}(A_d;m)
=
P\!\left(
\tilde q_{A_d}\ge \tilde q_{A_d}^{\mathrm{obs}}
\middle|\,
\mathrm{bkg}
\right),
\end{equation}
and
\begin{equation}
\mathrm{CL}_{s}(A_d;m)=
\frac{\mathrm{CL}_{s+b}(A_d;m)}{\mathrm{CL}_{b}(A_d;m)}.
\label{eq:cls-definition}
\end{equation}
Dividing by $\mathrm{CL}_b$ prevents a downward background fluctuation from producing
an artificially aggressive exclusion in a region with little intrinsic sensitivity
\cite{Junk1999,Read2002}. At confidence level $1-\alpha$, the single-dataset upper
limit is
\begin{equation}
A_{\mathrm{UL},d}(m;\alpha)
=\min\{A_d:\mathrm{CL}_s(A_d;m)\le\alpha\}.
\label{eq:aul-single}
\end{equation}
The analysis uses $\alpha=0.10$, corresponding to a 90\% CL upper limit; comparison
studies at 95\% use $\alpha=0.05$. We report the limit both in
signal-yield space and, after applying the dataset-specific conversion factor of
Eq.~\eqref{eq:Kd-conversion}, in $\eps^2$ space.
These are fixed-mass statements. The correlation-aware global-significance
calculation is introduced separately in Section~\ref{sec:v5-global}; it does
not change the pointwise $CL_s$ upper-limit definition.


\subsection{Asymptotic and toy calibration}

The local-significance mapping introduced above is treated asymptotically throughout
the scan. For upper limits, the same bounded statistic
$\tilde q_{A_d}$ is calibrated in one of two ways. In asymptotic mode, the
background-only Asimov dataset is used to evaluate the positive reference quantity
$q_A\equiv\tilde q_{A_d,A}>0$ for the tested nonzero amplitude. Writing
$q_{\mathrm{obs}}\equiv\tilde q_{A_d}^{\mathrm{obs}}$, the normal deviates entering
the two tail areas are
\begin{equation}
z_{s+b}=
\begin{cases}
\sqrt{q_{\mathrm{obs}}}, & q_{\mathrm{obs}}\le q_A,\\[2pt]
\dfrac{q_{\mathrm{obs}}+q_A}{2\sqrt{q_A}}, & q_{\mathrm{obs}}>q_A,
\end{cases}
\qquad
z_b=
\begin{cases}
\sqrt{q_A}-\sqrt{q_{\mathrm{obs}}}, & q_{\mathrm{obs}}\le q_A,\\[2pt]
\dfrac{q_A-q_{\mathrm{obs}}}{2\sqrt{q_A}}, & q_{\mathrm{obs}}>q_A.
\end{cases}
\label{eq:cls-asymptotic}
\end{equation}
The corresponding tail areas are
$\mathrm{CL}_{s+b}=1-\Phi(z_{s+b})$ and $\mathrm{CL}_b=\Phi(z_b)$, with
$\mathrm{CL}_s=\mathrm{CL}_{s+b}/\mathrm{CL}_b$.
At $q_{\rm obs}=0$ the statistic has a point mass, so inclusive toy tails
and the continuous asymptotic endpoint need a declared tie convention.
The implementation retains the continuous endpoint of the displayed
mapping, $\mathrm{CL}_{s+b}=1/2$ and $\mathrm{CL}_b=\Phi(\sqrt{q_A})$;
inclusive empirical tails at zero would both equal one. The reported
upper-limit roots lie on the positive-statistic branch.

The name \emph{Wald approximation} refers to treating the profiled objective
as locally quadratic, with an approximately normal unconstrained estimator:
\begin{equation}
 2\bigl[\ell_p(A)-\ell_p(\widehat A^{\rm unc})\bigr]
 \simeq\frac{(A-\widehat A^{\rm unc})^2}{\sigma_A^2}.
 \label{eq:v501-wald}
\end{equation}
Both branches of Eq.~\eqref{eq:cls-asymptotic} use this asymptotic idea. The
second branch accounts for the physical boundary when the unconstrained
best fit is negative. Its reference denominator is then the fit at zero,
so $\tilde q_A\simeq(A^2-2A\widehat A^{\rm unc})/\sigma_A^2$, rather than
$(A-\widehat A^{\rm unc})^2/\sigma_A^2$. This corresponds to
$q_{\rm obs}>q_A$ in the approximation. The implementation evaluates the
observed and Asimov profile statistics numerically and then applies this
piecewise probability mapping; it does not replace the underlying Poisson
fit with least squares. The approximation itself is not a coverage test.

In toy mode, the same statistic $\tilde q_{A_d}$ is calibrated directly with
background-only and signal-plus-background pseudoexperiments rather than by switching
to a different ordering variable. The shared-coupling combination described below uses
the identical logic after replacing the single-dataset amplitude $A_d$ with the common
parameter of interest $\eps^2$.

The pseudoexperiment program remains important even with the asymptotic formulae in
place because the HPS likelihood combines a multi-bin Poisson term, a GP-conditioned
blind-window covariance, a non-renormalized Gaussian signal template, and a
mass-dependent yield-to-$\eps^2$ conversion. The pull, spurious-signal, and
injection-recovery diagnostics of Section~\ref{sec:toys} test this machinery under the
declared conditional backgrounds. They are not direct \CLs{} coverage studies or a
scan-wide calibration of $Z$.

\subsection{Matched-reference injection--extraction diagnostics}
\label{sec:matched-reference-diagnostics}

We define the injection scale with a background-only reference fit that uses the same
support, binning, exclusion mask, kernel bounds, preprocessing, and GP refit as the
corresponding pseudoexperiment. If
$\mathcal{F}[\bm n;\bm w]$ denotes that full procedure for counts $\bm n$ and the
blind-window slice of the unit-normalized full-range signal template $\bm w$, then
\begin{equation}
(\widehat A_{\mathrm{ref}},\sigma_{A,\mathrm{ref}})
=\mathcal{F}[\bm n_{\mathrm{ref}};\bm w],
\qquad
A_{\mathrm{inj}}(z)=z\,\sigma_{A,\mathrm{ref}},
\qquad z\in\{0,1,3,5\}.
\label{eq:v4p9-matched-reference-axis}
\end{equation}
The reference uncertainty is a fixed calibration of the injection coordinate at a
given source, toy index, and mass.  It is not the post-injection uncertainty returned
for an individual extraction.

For pseudoexperiment $t$, the signed pull and paired response are
\begin{align}
p_t(z)&=\frac{\widehat A_t(z)-A_{\mathrm{inj},t}(z)}
                  {\sigma_{A,t}(z)},
\label{eq:v4p9-pull}\\
R_t(z)&=\frac{\widehat A_t(z)-\widehat A_t(0)}
                  {A_{\mathrm{inj},t}(z)},\qquad z>0.
\label{eq:v4p9-paired-response}
\end{align}
The zero-injection mean pull diagnoses a conditional spurious-signal offset under the
declared generating truth.  Its width probes the returned uncertainty scale, while
$R_t$ separates injection response from a background-specific zero-signal offset.
These quantities use the toy-level $\sigma_{A,t}$ in the denominator.  A finite
25-background sample supports two-sided 90\% Student-$t$ intervals on the mean and descriptive
widths, not a direct coverage statement.

Full refitting and fixed-background displays answer different questions. Holding the GP
fixed isolates the local profile-likelihood extraction.  Refitting the GP on every
pseudo-spectrum additionally tests background adaptation and possible signal
absorption. We use the latter matched-refit procedure. Numerical acceptance is decided
from likelihood, reproducibility, covariance, and kernel-coordinate diagnostics before
amplitudes or pulls are summarized.

\subsection{Observed-limit toy diagnostics}

The analysis uses two distinct families of $p$-values for different purposes. The
first is the discovery statistic of
Eq.~\eqref{eq:q0-single}, summarized through $q_0(m)$, $p_0(m)$, and
$Z_{\mathrm{local}}(m)$. The second consists of diagnostics
derived from the background-only \CLs{} limit ensemble at each mass. For
$N_{\mathrm{toys}}$ background-only pseudoexperiments and an observed limit
$\eps^2_{\mathrm{UL},\mathrm{obs}}(m)$ at the confidence level under study, those
quantities are
\begin{align}
p_{\mathrm{strong}}(m)
&=
\frac{1}{N_{\mathrm{toys}}}
\sum_{t=1}^{N_{\mathrm{toys}}}
\mathbf{1}\!\left[
\eps^2_{\mathrm{UL},t}(m)\le \eps^2_{\mathrm{UL},\mathrm{obs}}(m)
\right], \\
p_{\mathrm{weak}}(m)
&=
\frac{1}{N_{\mathrm{toys}}}
\sum_{t=1}^{N_{\mathrm{toys}}}
\mathbf{1}\!\left[
\eps^2_{\mathrm{UL},t}(m)\ge \eps^2_{\mathrm{UL},\mathrm{obs}}(m)
\right], \\
p_{\mathrm{two}}(m)
&=
\min\!\left[1,\,
2\min\!\left(p_{\mathrm{strong}}(m),\,p_{\mathrm{weak}}(m)\right)\right].
\end{align}
Here $p_{\mathrm{strong}}$ quantifies how unusually \emph{strong} the observed limit
is relative to the background-only toy ensemble, while $p_{\mathrm{weak}}$ quantifies
how unusually \emph{weak} it is. The two-sided quantity $p_{\mathrm{two}}$ is retained
only as a symmetric diagnostic; because it responds to both unusually strong
and unusually weak limits, it is less targeted than the one-sided discovery or
exclusion tests and is not used as a search statistic. Figure~\ref{fig:pvalue-tail-schematic}
gives the compact fixed-mass definition of these quantities, while
Figure~\ref{fig:pvalue-tail-examples} shows representative strong-limit, weak-limit,
and bounded median-case examples from the same toy-limit logic.

\begin{figure}[H]
\centering
\graphicorplaceholder{0.96\linewidth}{methodology_figs/pvalue_tail_schematic.png}
\caption{Illustrative toy-limit diagnostic at fixed mass. The left panel shows a
background-only ensemble of observed \CLs{} upper limits, and the right panel shows the
corresponding empirical cumulative distribution. The vertical line marks the observed
limit. The lower-tail area defines $p_{\mathrm{strong}}$, the upper-tail area defines
$p_{\mathrm{weak}}$, and
$p_{\mathrm{two}}=\min[1,2\min(p_{\mathrm{strong}},p_{\mathrm{weak}})]$ is used only
as a symmetric diagnostic summary.}
\label{fig:pvalue-tail-schematic}
\end{figure}

\begin{figure}[H]
\centering
\graphicorplaceholder{0.98\linewidth}{methodology_figs/pvalue_tail_examples.png}
\caption{Representative fixed-mass observed-limit toy diagnostics. The left panel
shows an unusually strong observed limit relative to the background-only toy
ensemble, the middle panel shows an unusually weak observed limit, and the right
panel shows a near-median case where the bounded symmetric summary saturates at
$p_{\mathrm{two}}=1$. These are limit diagnostics, not discovery
statistics.}
\label{fig:pvalue-tail-examples}
\end{figure}

These quantities are therefore limit diagnostics rather than discovery-significance
measures in the sense of $p_0$. If the observed limit lies near the median of the
background-only ensemble, then roughly half of the toys lie above it and half lie
below it, making both $p_{\mathrm{strong}}$ and $p_{\mathrm{weak}}$ naturally
$\mathcal{O}(0.5)$. Their finite resolution is $1/N_{\mathrm{toys}}$. In the
three-campaign combined ensemble, $N_{\mathrm{toys}}=300$, so a one-count change is
$1/300\simeq0.0033$ and a zero-count tail is reported as 0 of 300 rather than as an
exact zero probability. These toy-limit diagnostics cannot by themselves establish
or refute a high-significance local excess in the discovery statistic defined by the
one-sided local $p_0$ mapping above.

\subsection{Shared-\texorpdfstring{$\eps^2$}{epsilon2} combination across datasets}
\label{sec:combination}

A dark photon has one coupling, while each campaign has its own luminosity,
selected prompt rate and mass resolution. At a fixed mass, a common
$\psi\equiv\eps^2$ therefore predicts different signal yields in the active
datasets, as shown in Figure~\ref{fig:v501-common-coupling}.

\methodfig{0.99}{../figures/v501_common_coupling_flow.pdf}{One coupling predicts a
campaign-specific signal through its normalization and resolution. Each
campaign retains its own counts and correlated background constraint; the
likelihoods are multiplied and fitted with one common coupling. Only datasets
whose declared search contains the tested mass contribute.}
{fig:v501-common-coupling}

Let $\mathcal D(m)$ be that active set. We retain the signal template
$\bm w_d$, count mean $\bm b_d$ and covariance factor $L_d$ defined above.
In the electron-channel convention of Eq.~\eqref{eq:Kd-conversion}, the
radiative-fraction penalty gives
\begin{align}
 K_d^{\rm eff}(m)&=(1-\delta_{f,d})K_d(m),
 \label{eq:kdeff}\\
 \bm s_{{\rm unit},d}(m)&=K_d^{\rm eff}(m)\bm w_d(m),\\
 \bm\lambda_d(\psi,\bm\theta_d)
 &=\bm b_d+L_d\bm\theta_d+\psi\bm s_{{\rm unit},d}.
 \label{eq:v501-combined-mean}
\end{align}
The visible-branching correction used above the dimuon threshold is stated
separately below. The campaign likelihood and its product are
\begin{align}
 \mathcal L_d(\psi,\bm\theta_d)
 &=\left[\prod_i\operatorname{Pois}(n_{d,i}\mid\lambda_{d,i})\right]
   \exp\!\left(-\frac12\bm\theta_d^\top\bm\theta_d\right),
 \label{eq:dataset-poi-like}\\
 \mathcal L_{\rm comb}(\psi,\{\bm\theta_d\})
 &=\prod_{d\in\mathcal D(m)}\mathcal L_d(\psi,\bm\theta_d).
 \label{eq:combinedlike-product}
\end{align}
The background nuisance vectors are independent between campaigns. No
shared luminosity, radiative-fraction or resolution nuisance is profiled
in this baseline combination. This is the usual common-parameter structure
of a multi-channel likelihood~\cite{HistFactory,ATLASCMSHiggs2016}.

The implementation concatenates the count, background and unit-signal
vectors and uses a block-diagonal covariance. In compact notation,
\begin{align}
 \bm n&=\operatorname{col}_{d\in\mathcal D}\bm n_d,
 &\bm b&=\operatorname{col}_{d\in\mathcal D}\bm b_d,\\
 \bm s_{\rm unit}&=\operatorname{col}_{d\in\mathcal D}\bm s_{{\rm unit},d},
 &\bm\theta&=\operatorname{col}_{d\in\mathcal D}\bm\theta_d,\\
 L_{\rm comb}&=\operatorname{blockdiag}_{d\in\mathcal D}L_d,
 &\bm C_{\rm comb}&=L_{\rm comb}L_{\rm comb}^\top.
\end{align}
Here $\operatorname{col}$ stacks column vectors without combining their
bins. The single vector mean is
\begin{equation}
 \bm\lambda(\psi,\bm\theta)
 =\bm b+L_{\rm comb}\bm\theta+\psi\bm s_{\rm unit}.
\end{equation}
Thus all campaigns are fitted at the same $\psi$, while the nuisance blocks
respond separately to their spectra.

It is the \emph{profiled negative log likelihoods} that add at fixed coupling:
\begin{align}
 \ell_{p,\rm comb}(\psi)
 &=\sum_{d\in\mathcal D(m)}\ell_{p,d}(\psi),
 \label{eq:v501-common-profile-sum}\\
 \widehat\psi^{\rm unc}
 &=\operatorname*{arg\,min}_{\psi\in\mathbb R}\ell_{p,\rm comb}(\psi),
 &\widehat\psi^+&=\max(0,\widehat\psi^{\rm unc}).
\end{align}
The auxiliary signed fit is used for extraction diagnostics. For physical
upper limits the bounded statistic is
\begin{equation}
 \widetilde q_{\psi}=
 \begin{cases}
 0,&\widehat\psi^{\rm unc}>\psi,\\
 2\displaystyle\sum_{d\in\mathcal D(m)}
 \bigl[\ell_{p,d}(\psi)-\ell_{p,d}(\widehat\psi^+)\bigr],
 &\widehat\psi^{\rm unc}\leq\psi,
 \end{cases}
 \qquad \psi\geq0.
 \label{eq:v501-common-bounded-statistic}
\end{equation}
Every denominator is evaluated at the \emph{same common best fit}
$\widehat\psi^+$. Adding likelihood ratios normalized to separate
campaign best fits would define a different statistic. Likewise, neither
averaging upper limits nor adding local significances performs this fit.
The shared-coupling upper limit satisfies
\begin{equation}
 \eps^2_{\rm UL}(m;\alpha)
 =\min\{\psi\geq0:\mathrm{CL}_s(\psi;m)\leq\alpha\},
 \qquad \alpha=0.10.
\end{equation}
The combined local excess test uses this same profiled objective at zero
and at $\widehat\psi^+$. Its gain comes from the information supplied by
independent campaigns about one physical parameter, while retaining each
campaign's normalization and background uncertainty.

\subsection{Global significance and the look-elsewhere treatment}

The most extreme local excess in a declared scan is summarized by its minimum
one-sided $p_0$ or maximum $q_0$,
\begin{equation}
p_{0,\min}^{\mathrm{obs}} = \min_m p_0^{\mathrm{obs}}(m),
\qquad
q_{0,\max}^{\mathrm{obs}} = \max_m q_0^{\mathrm{obs}}(m).
\end{equation}
The cleanest global-significance definition is then scan based:
\begin{equation}
p_{\mathrm{global}}^{\mathrm{toy}}
=
\frac{1}{N_{\mathrm{toys}}}
\sum_{t=1}^{N_{\mathrm{toys}}}
\mathbf{1}\!\left[
q_{0,\max}^{(t)} \ge q_{0,\max}^{\mathrm{obs}}
\right],
\label{eq:pglobal-toy}
\end{equation}
with
\begin{equation}
Z_{\mathrm{global}}^{\mathrm{toy}} = \Phi^{-1}\!\left(1-p_{\mathrm{global}}^{\mathrm{toy}}\right).
\end{equation}
The normal-quantile conversion requires a resolved tail probability. With
zero exceedances, the result is the one-sided 95\% sampling bound
$p_{\rm global}<1-0.05^{1/N_{\rm toys}}$; an infinite significance is not
assigned. Section~\ref{sec:v5-global} applies this convention to the
conditional full-field calculation.
Earlier versions used an effective-trials approximation based on the scan span
and detector resolution as an analytic look-elsewhere reference \cite{GrossVitells2010}:
\begin{equation}
p_{\mathrm{global}}^{\mathrm{Sidak}}
=
1-\left(1-p_{0,\min}^{\mathrm{obs}}\right)^{N_{\mathrm{eff}}},
\qquad
Z_{\mathrm{global}}^{\mathrm{Sidak}}
=
\Phi^{-1}\!\left(1-p_{\mathrm{global}}^{\mathrm{Sidak}}\right).
\label{eq:sidak-global}
\end{equation}
For small $p_{0,\min}^{\mathrm{obs}}$, Eq.~\eqref{eq:sidak-global} reduces to the
familiar $p_{\mathrm{global}}\simeq N_{\mathrm{eff}}p_{0,\min}$ form.

This effective-trials expression is retained as a historical heuristic. A
scan span and resolution scale do not by themselves model the nonstationary
correlations caused by the moving GP exclusion or changing campaign membership.
The conditional full-field calculation in Section~\ref{sec:v5-global} evaluates
those correlations explicitly. Neither the heuristic nor additional Gaussian
field sampling replaces validation of the generating background and the
accessible tails with coherent complete Poisson scans.

\subsection{Signal absorption and alternative background-profile treatments}

We describe blind-window background uncertainty with additive Gaussian-constrained
nuisance parameters derived from the GP mean and covariance. Normalizing the signal
over its full line shape removes the bias from blind-window renormalization, but GP
refitting or the nuisance model can still absorb part of an injected signal.

The most likely residual sources are the refitting freedom of the GP hyperparameters,
especially the effective length scale, and the flexibility of the Gaussian-profile
background nuisance itself. We therefore retain fixed-background and alternative
Poisson treatments as targeted robustness tests if matched pseudoexperiments reveal
unacceptable absorption.

\subsection{Dimuon-threshold interpretation}
\label{sec:methodology-dimuon-threshold}

The yield-to-coupling conversion in Eq.~\eqref{eq:eps2-conversion} is an
electron-channel normalization: the fitted resonance amplitude is the
$A^\prime\to e^+e^-$ yield visible in the prompt invariant-mass spectrum. For the
minimal kinetically mixed dark photon, this is also the full visible branching
interpretation below the dimuon threshold, $m_{A^\prime}<2m_\mu=\SI{211.3167}{MeV}$,
where the $\mu^+\mu^-$ final state is kinematically closed. This is the same
fixed-target convention used in the published HPS prompt searches
\cite{EssigEtAl2011,HPS2015DarkPhoton,HPS2016PromptLong}. Above threshold, the same
electron-channel yield corresponds to a weaker minimal-model limit on $\eps^2$ because
the total width includes the newly open $A^\prime\to\mu^+\mu^-$ mode
\cite{Holdom1986,Fabbrichesi2020DarkPhoton}.

Neglecting the electron mass and remaining below the two-pion hadronic threshold, the
charged-lepton width ratio is
\begin{equation}
R_{\mu/e}(m_{A^\prime})
\equiv
\frac{\Gamma(A^\prime\to\mu^+\mu^-)}
{\Gamma(A^\prime\to e^+e^-)}
=
\sqrt{1-\frac{4m_\mu^2}{m_{A^\prime}^2}}\,
\left(1+\frac{2m_\mu^2}{m_{A^\prime}^2}\right),
\qquad m_{A^\prime}>2m_\mu,
\label{eq:dimuon-width-ratio}
\end{equation}
with $R_{\mu/e}=0$ below threshold. The inverse electron branching fraction used for
the minimal-model reinterpretation is therefore
\begin{equation}
N_{\mathrm{eff}}^{\mathrm{BR}}(m_{A^\prime})
\equiv
\frac{1}{\mathcal{B}(A^\prime\to e^+e^-)}
=
1+R_{\mu/e}(m_{A^\prime}),
\label{eq:dimuon-neff}
\end{equation}
and an electron-channel coupling limit is mapped to the minimal visible-dark-photon
convention by
\begin{equation}
\eps^2_{\mathrm{min}}(m_{A^\prime})
=
N_{\mathrm{eff}}^{\mathrm{BR}}(m_{A^\prime})\,\eps^2_{ee}(m_{A^\prime}).
\label{eq:dimuon-eps2-map}
\end{equation}
This operation does not change the local fit, the signal-yield upper limit, the
\CLs{} value, or the local $p_0$ scan; it changes only the coupling interpretation of
an electron-channel result. Figure~\ref{fig:methodology-dimuon-factor} shows the
factor applied in the 90\% CL result obtained with \CLs{}. The factor is unity below
$2m_\mu$, rises continuously above threshold, and reaches
$N_{\mathrm{eff}}^{\mathrm{BR}}=1.725$ at \SI{250}{MeV}
($\mathcal{B}_{ee}=0.580$).

The mass range of a scan-global discovery statement must equal its declared
search range~\cite{GrossVitells2010}. A branching-fraction conversion changes the
physical interpretation of the coupling, not that search definition. Pointwise
\CLs{} exclusion curves are not multiplied by a look-elsewhere trials factor.

\begin{figure}[H]
\centering
\graphicorplaceholder{0.78\linewidth}{final_limit_projection_figs/90cls/dimuon/dimuon_factor_vs_mass.png}
\caption{Minimal-model branching factor used to reinterpret electron-channel
$\eps^2$ limits above the dimuon threshold. The red curve shows
$N_{\mathrm{eff}}^{\mathrm{BR}}=1/\mathcal{B}(A^\prime\to e^+e^-)$ and the blue dashed
curve shows $\mathcal{B}(A^\prime\to e^+e^-)$.}
\label{fig:methodology-dimuon-factor}
\end{figure}

With the observable, background prediction, signal model, and likelihood specified,
the remaining question is empirical: does this procedure create a false narrow signal,
misstate its uncertainty, or absorb a signal that is actually present? We answer that
question with pseudoexperiments in the next section.
